\section{Cauchy Principal Value}

\begin{enumerate}
  \item What problem is the Cauchy principal value trying to solve?
  \item Why does
    \[
      \int_{-1}^{1}\frac{1}{x},dx
    \]
    fail to exist as an ordinary improper integral?
  \item What is the difference between an improper integral and a principal value?
  \item Why is a principal value not simply another way to evaluate an improper integral?
  \item What does the word ``symmetric'' mean in the definition of principal value?
  \item Why is symmetry essential in the definition?
  \item Can a principal value exist when the corresponding improper integral does not?
  \item Can an improper integral exist when the principal value also exists?
  \item What information is lost when we replace an improper integral by a principal value?
  \item What geometric intuition can be attached to principal values?
\end{enumerate}

\subsection{Computational Questions}

\begin{enumerate}[resume]
  \item Compute
    \[
      \operatorname{PV}\int_{-1}^{1}\frac{1}{x},dx.
    \]
  \item Why does the answer not contradict the fact that the ordinary improper integral diverges?
  \item Compute
    \[
      \operatorname{PV}\int_{-2}^{3}\frac{1}{x},dx.
    \]
  \item Why is the answer no longer zero?
  \item Derive a general formula for
    \[
      \operatorname{PV}\int_a^b\frac{1}{x},dx,
    \]
    assuming \(a<0<b\).
  \item Compute
    \[
      \operatorname{PV}\int_{-1}^{1}\frac{x}{x^2+1},dx.
    \]
  \item Compute
    \[
      \operatorname{PV}\int_{-1}^{1}\frac{1+x}{x},dx.
    \]
  \item Why is it useful to split the integrand into a singular part and a regular part?
  \item Compute
    \[
      \operatorname{PV}\int_{-1}^{1}\frac{\sin x}{x},dx.
    \]
  \item Is a principal value always finite when the singularity is of the form \(1/x\)?
\end{enumerate}

\subsection{Questions About Singularities}

\begin{enumerate}[resume]
  \item Why does \(1/x\) behave differently from \(1/x^2\)?
  \item Does
    \[
      \operatorname{PV}\int_{-1}^{1}\frac{1}{x^2},dx
    \]
    exist?
  \item If not, why not?
  \item What property of \(1/x\) allows cancellation?
  \item What property of \(1/x^2\) prevents cancellation?
  \item For which powers \(p\) does
    \[
      \operatorname{PV}\int_{-1}^{1}\frac{1}{x^p},dx
    \]
    exist?
  \item What role does oddness or evenness of a function play?
  \item When does symmetry force a principal value to be zero?
  \item Can an odd function have a divergent principal value?
  \item Can an even function have a finite principal value?
\end{enumerate}

\subsection{Questions About Infinity}

\begin{enumerate}[resume]
  \item What is
    \[
      \operatorname{PV}\int_{-\infty}^{\infty}\frac{1}{x},dx?
    \]
  \item How does this differ from
    \[
      \int_{-\infty}^{\infty}\frac{1}{x},dx?
    \]
  \item Why must limits at \(+\infty\) and \(-\infty\) be taken symmetrically?
  \item What happens if the two limits are taken independently?
  \item Give an example where symmetric truncation gives a finite value but ordinary convergence fails.
  \item Why is
    \[
      \operatorname{PV}\int_{-\infty}^{\infty}\sin x,dx
    \]
    an interesting example?
  \item How does principal value at infinity differ from principal value around an interior singularity?
\end{enumerate}

\subsection{Theory Questions}

\begin{enumerate}[resume]
  \item State the formal definition of the Cauchy principal value for an interior singularity.
  \item State the formal definition for an integral over the whole real line.
  \item Why is the principal value defined as a limit?
  \item Under what conditions does the principal value agree with the ordinary improper integral?
  \item Is the principal value linear?
  \item When is it legitimate to split a principal value integral into two pieces?
  \item Can principal values always be manipulated algebraically as ordinary integrals?
  \item What pitfalls arise when changing variables in principal value integrals?
  \item Why must one be careful about symmetry when performing substitutions?
  \item Is principal value integration compatible with all standard integration theorems?
  \item Which familiar theorems from elementary calculus continue to work unchanged?
  \item Which require extra hypotheses?
  \item Why do analysts treat principal values as generalized notions of integration?
\end{enumerate}

\end{document}
